% segment-local-id: OR03-S006
% segment-id: d90.orig.v1.tr03.seg0006
% license: CC BY-SA 4.0
% provenance: Materi asesmen asli berbahasa Indonesia, disusun secara independen.

\section{Set soal IV: metode stokastik}
\label{orig03:ps04}

% stable-id: d90.orig.v1.tr03.ps04.exercise.0001
\begin{exercise}[Pencuplikan tak seragam dan koreksi kepentingan]
\label{orig03:ps04:exercise:0001}
Misalkan
\[
  F(x)=\frac1N\sum_{i=1}^N f_i(x),
  \qquad g_i=\nabla f_i(x),
\]
dan, bersyarat pada informasi masa lalu $\mathcal F$, indeks $I$ dipilih dengan
peluang $p_i>0$, $\sum_i p_i=1$. Peluang dan titik $x$ boleh
$\mathcal F$-terukur.

\begin{enumerate}
  \item Hitung bias penduga naif $H=g_I$ terhadap $\nabla F(x)$.
  \item Tunjukkan bahwa $G=g_I/(Np_I)$ tidak bias secara bersyarat.
  \item Turunkan rumus eksak untuk
  $\mathbb E[\lVert G-\nabla F(x)\rVert^2\mid\mathcal F]$.
\end{enumerate}
\end{exercise}

% stable-id: d90.orig.v1.tr03.ps04.hint.0001
\paragraph{Petunjuk bertahap.}
\label{orig03:ps04:hint:0001}
\textbf{Tahap 1.} Tulis ekspektasi bersyarat sebagai jumlah atas semua nilai
yang mungkin bagi $I$.
\textbf{Tahap 2.} Setelah memperoleh ketakbiasan, gunakan identitas
$\mathbb E\lVert G-m\rVert^2=\mathbb E\lVert G\rVert^2-\lVert m\rVert^2$
dengan $m=\mathbb E[G\mid\mathcal F]$.

% stable-id: d90.orig.v1.tr03.ps04.answer.0001
\paragraph{Jawaban singkat.}
\label{orig03:ps04:answer:0001}
Penduga naif mempunyai bias
$\sum_i p_i g_i-N^{-1}\sum_i g_i$. Sementara itu,
\[
 \mathbb E[G\mid\mathcal F]=\nabla F(x),\qquad
 \mathbb E[\lVert G-\nabla F(x)\rVert^2\mid\mathcal F]
 =\frac1{N^2}\sum_{i=1}^N\frac{\lVert g_i\rVert^2}{p_i}
  -\lVert\nabla F(x)\rVert^2.
\]

% stable-id: d90.orig.v1.tr03.ps04.solution.0001
\paragraph{Solusi lengkap.}
\label{orig03:ps04:solution:0001}
Karena $p_i$ dan $g_i$ diketahui setelah mengondisikan pada $\mathcal F$,
\[
 \mathbb E[H\mid\mathcal F]=\sum_i p_i g_i.
\]
Jadi selisihnya dari gradien penuh adalah
$\sum_i(p_i-1/N)g_i$, yang umumnya tidak nol. Untuk penduga terkoreksi,
\[
 \mathbb E[G\mid\mathcal F]
 =\sum_i p_i\frac{g_i}{Np_i}
 =\frac1N\sum_i g_i=\nabla F(x).
\]
Selanjutnya,
\[
 \mathbb E[\lVert G\rVert^2\mid\mathcal F]
 =\sum_i p_i\frac{\lVert g_i\rVert^2}{N^2p_i^2}
 =\frac1{N^2}\sum_i\frac{\lVert g_i\rVert^2}{p_i}.
\]
Mengurangkan kuadrat norma mean bersyarat memberi rumus varians yang
dinyatakan pada jawaban singkat. Positivitas setiap $p_i$ penting agar koreksi
terdefinisi dan tidak menghilangkan komponen mana pun.

% stable-id: d90.orig.v1.tr03.ps04.exercise.0002
\begin{exercise}[Ketaksamaan tiga titik untuk langkah Bregman komposit]
\label{orig03:ps04:exercise:0002}
Pada himpunan konveks tertutup $C$, misalkan $h$ terdiferensialkan dan konveks,
$r$ konveks proper semikontinu-bawah, dan
\[
 x^+\in\arg\min_{x\in C}
 \left\{\eta\langle g,x\rangle+\eta r(x)+D_h(x,x_k)\right\},
 \qquad \eta>0.
\]
Buktikan bahwa untuk setiap $u\in C$,
\[
 \eta\bigl(\langle g,x^+-u\rangle+r(x^+)-r(u)\bigr)
 \le D_h(u,x_k)-D_h(u,x^+)-D_h(x^+,x_k).
\]
Nyatakan dengan jelas kondisi optimalitas yang Anda gunakan.
\end{exercise}

% stable-id: d90.orig.v1.tr03.ps04.hint.0002
\paragraph{Petunjuk bertahap.}
\label{orig03:ps04:hint:0002}
\textbf{Tahap 1.} Pilih $s^+\in\partial r(x^+)$ dan
$n^+\in N_C(x^+)$ dari inklusi optimalitas.
\textbf{Tahap 2.} Gunakan ketaksamaan subgradien dan identitas tiga titik
untuk divergensi Bregman.

% stable-id: d90.orig.v1.tr03.ps04.answer.0002
\paragraph{Jawaban singkat.}
\label{orig03:ps04:answer:0002}
Inklusi
\[
 0\in\eta g+\eta\partial r(x^+)+\nabla h(x^+)-\nabla h(x_k)+N_C(x^+)
\]
bersama dengan
\[
 \langle\nabla h(x^+)-\nabla h(x_k),u-x^+\rangle
 =D_h(u,x_k)-D_h(u,x^+)-D_h(x^+,x_k)
\]
memberikan hasilnya.

% stable-id: d90.orig.v1.tr03.ps04.solution.0002
\paragraph{Solusi lengkap.}
\label{orig03:ps04:solution:0002}
Ambil $s^+\in\partial r(x^+)$ dan $n^+\in N_C(x^+)$ sehingga
\[
 0=\eta g+\eta s^++\nabla h(x^+)-\nabla h(x_k)+n^+.
\]
Karena $u\in C$, definisi kerucut normal memberi
$\langle n^+,u-x^+\rangle\le0$. Mengalikan persamaan optimalitas dengan
$x^+-u$ menghasilkan
\[
 \eta\langle g,x^+-u\rangle+\eta\langle s^+,x^+-u\rangle
 \le \langle\nabla h(x^+)-\nabla h(x_k),u-x^+\rangle.
\]
Ketaksamaan subgradien menyatakan
$\langle s^+,x^+-u\rangle\ge r(x^+)-r(u)$. Ruas kanan sama dengan selisih
tiga divergensi Bregman pada jawaban singkat. Substitusi kedua fakta itu
menyelesaikan bukti.

% stable-id: d90.orig.v1.tr03.ps04.exercise.0003
\begin{exercise}[Varians terbatas bukan momen kedua seragam]
\label{orig03:ps04:exercise:0003}
Pertimbangkan $f(x)=x^2/2$ dan orakel stokastik
$G(x,\xi)=x+\xi$, dengan $\mathbb E\xi=0$ dan
$\mathbb E\xi^2=\sigma^2<\infty$.

\begin{enumerate}
  \item Verifikasi ketakbiasan dan varians bersyarat yang seragam.
  \item Tunjukkan bahwa tidak ada batas seragam untuk
  $\mathbb E|G(x,\xi)|^2$ ketika $x$ menjelajahi $\mathbb R$.
  \item Buktikan arah sebaliknya: jika
  $\mathbb E\lVert G(x,\xi)\rVert^2\le M^2$ dan $G$ tidak bias, maka
  $\lVert\nabla f(x)\rVert\le M$.
\end{enumerate}
\end{exercise}

% stable-id: d90.orig.v1.tr03.ps04.hint.0003
\paragraph{Petunjuk bertahap.}
\label{orig03:ps04:hint:0003}
\textbf{Tahap 1.} Pisahkan $G-\nabla f$ dari $G$ sendiri.
\textbf{Tahap 2.} Untuk arah sebaliknya, terapkan ketaksamaan Jensen pada
fungsi kuadrat norma.

% stable-id: d90.orig.v1.tr03.ps04.answer.0003
\paragraph{Jawaban singkat.}
\label{orig03:ps04:answer:0003}
$G$ tidak bias dan
$\mathbb E|G-\nabla f(x)|^2=\sigma^2$, tetapi
$\mathbb E|G|^2=x^2+\sigma^2$, sehingga tak terbatas seragam. Sebaliknya,
$\lVert\nabla f(x)\rVert^2=\lVert\mathbb E G\rVert^2
\le\mathbb E\lVert G\rVert^2\le M^2$.

% stable-id: d90.orig.v1.tr03.ps04.solution.0003
\paragraph{Solusi lengkap.}
\label{orig03:ps04:solution:0003}
Gradien eksak ialah $\nabla f(x)=x$. Karena mean $\xi$ nol,
$\mathbb E[G(x,\xi)]=x$, sedangkan
$G(x,\xi)-\nabla f(x)=\xi$. Jadi varians galat orakel selalu
$\sigma^2$, tidak bergantung pada $x$. Namun,
\[
 \mathbb E|x+\xi|^2=x^2+2x\mathbb E\xi+\mathbb E\xi^2
 =x^2+\sigma^2,
\]
yang menuju tak hingga saat $|x|\to\infty$. Dengan demikian, asumsi varians
terbatas tidak boleh diam-diam diganti dengan asumsi momen kedua seragam.
Untuk arah terakhir, ketakbiasan dan konveksitas kuadrat norma memberi
\[
 \lVert\nabla f(x)\rVert^2
 =\lVert\mathbb E G(x,\xi)\rVert^2
 \le\mathbb E\lVert G(x,\xi)\rVert^2\le M^2.
\]

% stable-id: d90.orig.v1.tr03.ps04.exercise.0004
\begin{exercise}[Bias bersyarat pada pengacakan ulang]
\label{orig03:ps04:exercise:0004}
Dalam satu epok, ambil permutasi seragam $\pi$ dari
$\{1,\ldots,N\}$. Setelah indeks
$\pi_1,\ldots,\pi_t$ terungkap, misalkan $R_t$ adalah himpunan indeks yang
tersisa dan $\mathcal F_t$ memuat riwayat tersebut.

\begin{enumerate}
  \item Untuk titik beku $x$, hitung
  $\mathbb E[\nabla f_{\pi_{t+1}}(x)\mid\mathcal F_t]$.
  \item Tunjukkan bahwa jumlah gradien selama satu epok pada titik beku tepat
  sama dengan $N\nabla F(x)$.
  \item Jelaskan mengapa bukti SGD yang mengandalkan selisih martingal tidak
  dapat langsung dipakai ketika titik iterasi berubah di dalam epok.
\end{enumerate}
\end{exercise}

% stable-id: d90.orig.v1.tr03.ps04.hint.0004
\paragraph{Petunjuk bertahap.}
\label{orig03:ps04:hint:0004}
\textbf{Tahap 1.} Bersyarat pada riwayat, indeks berikutnya seragam hanya pada
$R_t$, bukan pada seluruh himpunan indeks.
\textbf{Tahap 2.} Bedakan identitas kombinatorial pada satu titik $x$ dari
jumlah $\nabla f_{\pi_t}(x_t)$ pada titik-titik yang berlainan.

% stable-id: d90.orig.v1.tr03.ps04.answer.0004
\paragraph{Jawaban singkat.}
\label{orig03:ps04:answer:0004}
\[
 \mathbb E[\nabla f_{\pi_{t+1}}(x)\mid\mathcal F_t]
 =\frac1{|R_t|}\sum_{i\in R_t}\nabla f_i(x),
 \qquad
 \sum_{t=1}^N\nabla f_{\pi_t}(x)=N\nabla F(x).
\]
Mean bersyarat pertama umumnya bukan gradien penuh; bila $x_t$ berubah,
identitas jumlah pada titik beku juga tidak berlaku.

% stable-id: d90.orig.v1.tr03.ps04.solution.0004
\paragraph{Solusi lengkap.}
\label{orig03:ps04:solution:0004}
Setelah $t$ indeks terungkap, setiap elemen $R_t$ mempunyai peluang
$1/|R_t|$ untuk muncul berikutnya. Karena itu mean bersyarat adalah rata-rata
gradien komponen yang tersisa. Rata-rata ini hanya sama dengan gradien penuh
dalam keadaan khusus. Di sisi lain, sebuah permutasi memuat setiap indeks
tepat sekali, sehingga untuk satu titik yang sama,
\[
 \sum_{t=1}^N\nabla f_{\pi_t}(x)=\sum_{i=1}^N\nabla f_i(x)=N\nabla F(x).
\]
Pada algoritme sebenarnya, suku ke-$t$ dinilai di $x_t$, sehingga jumlahnya
adalah $\sum_t\nabla f_{\pi_t}(x_t)$ dan tidak lagi dapat diganti oleh jumlah
gradien pada satu titik. Selain itu, galat gradien terhadap gradien penuh
umumnya mempunyai mean bersyarat tak nol. Maka argumen selisih martingal SGD
dengan pengambilan sampel independen memerlukan analisis baru, bukan sekadar
penggantian notasi.

% stable-id: d90.orig.v1.tr03.ps04.exercise.0005
\begin{exercise}[Lantai galat akibat derau persisten]
\label{orig03:ps04:exercise:0005}
Misalkan barisan tak negatif memenuhi
\[
 a_{k+1}\le(1-\mu\eta)a_k+\eta^2\sigma^2,
 \qquad 0<\mu\eta<1.
\]
Turunkan batas tertutup untuk $a_k$, hitung limit superiornya, dan jelaskan
apa yang harus diubah agar jaminan konvergensi eksak tetap masuk akal ketika
$\sigma>0$.
\end{exercise}

% stable-id: d90.orig.v1.tr03.ps04.hint.0005
\paragraph{Petunjuk bertahap.}
\label{orig03:ps04:hint:0005}
\textbf{Tahap 1.} Tetapkan $q=1-\mu\eta$ dan buka rekursi sebagai deret
geometrik.
\textbf{Tahap 2.} Bandingkan faktor kontraksi $q^k$ dengan suku paksa yang
tidak lenyap.

% stable-id: d90.orig.v1.tr03.ps04.answer.0005
\paragraph{Jawaban singkat.}
\label{orig03:ps04:answer:0005}
Dengan $q=1-\mu\eta$,
\[
 a_k\le q^k a_0+\frac{\eta\sigma^2}{\mu}(1-q^k),
 \qquad
 \limsup_{k\to\infty}a_k\le\frac{\eta\sigma^2}{\mu}.
\]
Langkah tetap dengan derau persisten hanya menjamin lingkungan solusi;
langkah menurun atau reduksi varians diperlukan untuk meniadakan lantai ini.

% stable-id: d90.orig.v1.tr03.ps04.solution.0005
\paragraph{Solusi lengkap.}
\label{orig03:ps04:solution:0005}
Iterasi rekursi memberikan
\[
 a_k\le q^k a_0+\eta^2\sigma^2\sum_{j=0}^{k-1}q^j
 =q^k a_0+\eta^2\sigma^2\frac{1-q^k}{1-q}.
\]
Karena $1-q=\mu\eta$, suku kedua sama dengan
$\eta\sigma^2(1-q^k)/\mu$. Hipotesis memastikan $0<q<1$, sehingga
$q^k\to0$ dan batas limit superior mengikuti. Bila $\eta$ tetap dan
$\sigma>0$, suku paksa tidak hilang; rekursi itu sendiri tidak dapat
membuktikan $a_k\to0$. Jadwal langkah yang menuju nol dengan syarat penjumlahan
yang sesuai, atau penduga dengan varians yang ikut menuju nol, mengubah
rekursi sehingga lantai galat dapat lenyap.
